\begin{abstract}

高比例光伏接入下，微网购电与储能调度受到能量跨期耦合、预测误差、信息分时到达及动态电价的共同影响。若仅依据点预测求解单一最优轨迹，预测偏差将通过高额紧急购电与计划调整成本非对称地传导至实际运行费用。本文构建基于渐进信息集的序贯优化框架，将确定性调度、两阶段随机规划与多时点滚动重优化统一于同一预测--决策经济评价体系中。

针对问题一，将储能的跨时段能量转移刻画为\textbf{确定性混合整数线性规划}问题，在满足功率平衡及储能运行边界的条件下求解最优购电策略。最优全天计划购电量为 \textbf{59,482.70 kWh}，购电费用为 \textbf{35,126.95 元}，较无储能基准降低 \textbf{26.90\%}。\textbf{LP松弛检验显示模型整数间隙为零，且未出现同时充放电现象。}

针对问题二，严格遵循\textbf{因果序贯评价协议}，采用周周期负荷预测与七日同刻光伏季节均值作为预测基线，并\textbf{基于历史负荷--光伏成对整日残差生成联合场景}，构建\textbf{两阶段随机MILP}确定日前购电计划，同时建立严格的物理执行边界与跨日储能推演。2025年2月至12月连续334天均获得最优解，\textbf{实际总费用为 14,839,462.00元，紧急购电量为 398,161.90 kWh}。结果表明，在5倍惩罚下，预测误差会非对称传导至成本，模型有效性应由下游调度经济性而非单一统计指标决定。

针对问题三，构建\textbf{提前期分段的融合预测与序贯滚动重优化模型}，并以\textbf{边际信息价值（VOI）}衡量新增预报的经济贡献。四组信息扩张策略中，\textbf{$S_2=\{0,6,12\}$ 的年度费用最低，为 $14,505,500.73$ 元}；\textbf{$6{:}00$ 和 $12{:}00$ 更新分别产生 $95,120.80$ 元和 $127,349.06$ 元的边际收益，而 $18{:}00$ 更新的实现边际价值为 $-2,217.41$ 元}。Bootstrap重采样检验（$95\%$ 置信区间 $[546.36,\,792.54]$ 元/天）进一步验证了$S_2$相较$S_0$（日均节省$666.08$元）具有显著且稳定的日级经济收益。

针对问题四，将附件4给出的10 min动态电价作为\textbf{外生时变参数嵌入上述随机--滚动框架}。动态价格下，\textbf{日前随机调度总费用为 15,547,956.36元，滚动策略 $S_2$ 的费用降至 15,196,922.47元}，较前者降低 2.26\%，紧急购电量下降 23.74\%。各预报节点的边际信息价值排序与固定电价情形保持一致，表明所建序贯决策框架对价格机制变化具有较好的稳定性。

综上，本文将四问统一为由\textbf{确定性最优--随机日前决策--序贯信息更新--动态价格响应}逐级扩展的运筹优化问题。研究表明，微网储能调度的关键不在于单纯追求更低的预测误差，而在于将预测不确定性、储能物理状态与市场结算规则共同转化为可执行的经济决策；预测信息的有效性最终取决于其能够贡献的\textbf{边际决策价值}。

\keywords{微网能量管理\quad 联合残差场景\quad 两阶段随机规划\quad 滚动重优化\quad 信息价值}

\end{abstract}
